% Standalone English source for Overleaf; no local LaTeX compilation.
% The new statements in this note are paper proofs, not new Lean theorems.
\documentclass[11pt]{article}
\usepackage[T1]{fontenc}
\usepackage[utf8]{inputenc}
\usepackage{lmodern}
\usepackage[a4paper,margin=25mm]{geometry}
\usepackage{amsmath,amssymb,amsthm,mathtools}
\usepackage{booktabs,longtable,array}
\usepackage{enumitem}
\usepackage{xcolor}
\usepackage{hyperref}
\usepackage{xurl}
\hypersetup{colorlinks=true,linkcolor=blue!55!black,urlcolor=blue!55!black}
\setlength{\emergencystretch}{3em}
\setlist{itemsep=3pt,topsep=5pt}
\newtheorem{theorem}{Theorem}[section]
\newtheorem{lemma}[theorem]{Lemma}
\newtheorem{definition}[theorem]{Definition}
\theoremstyle{remark}
\newtheorem{remark}[theorem]{Remark}
\newcommand{\N}{\mathbb N}
\newcommand{\Q}{\mathbb Q}
\newcommand{\EF}{\operatorname{ExactFeasible}}
\newcommand{\code}[1]{\nolinkurl{#1}}
\newcommand{\src}[2]{\href{https://github.com/huam54925-cpu/uav-routing-certification/blob/1bc992475a9bd515d77245e2bc0471e5ecab2fd5/LeanCompletion/UAV/#1.lean}{#2}}
\title{Allocation Boundaries in Exact Tree Routing\\
  \large A Strict-Hub Three-Partition Correspondence and an Internal-Subtree Gluing Test}
\author{Project Research Note}
\date{6 November 2026}
\begin{document}
\maketitle

\begin{abstract}
This note supplies two bounded paper results requested alongside the first representation-bridge package.
For Issue \#23 it proves that a specified strict-hub exact-routing instance is feasible if and only if the source item instance has a three-partition.
For Issue \#19 it gives two explicit partial route families with the same aggregate subtree flow and identical per-route boundary crossing counts, but different compatibility with one fixed exterior context.
Both arguments preserve per-occurrence service charges and the prohibition on intermediate depot visits.
They are mathematical proofs, not new Lean results or complexity-classification claims.
\end{abstract}

\section{Common Routing Semantics}

A legal route has a nonempty task sequence, starts and ends at the depot $o$, has no intermediate occurrence of $o$, and has distinct consecutive vertices.
Repeated non-depot visits are allowed and each occurrence is charged its service fee.
For a route $r$, let $F_r(u,v)$ be its directed edge multiplicity, $N_r(v)$ its task-visit count, and
\[
 C_s(r)=\sum_v s(v)N_r(v).
\]
An exact feasible family of size $m$ has exactly the prescribed total visit counts at non-depot vertices, exactly the prescribed total directed flow $x$, and $C_s(r)\leq T$ for every route.
These are the existing \src{RouteToFlow}{\code{LegalRoute}} and
\src{SupportExactFeasibility}{\code{ExactFeasible}} semantics.
Graph support is imposed separately; it is not a field of \code{LegalRoute}.

For a rooted tree $R$ and nonroot parent-edge multiplicities $H$, let $E_R(H)$ use each parent-child edge $H(v)$ times in each direction.
The required non-depot visit count is
\begin{equation}\label{eq:demand}
 d(v)=H(v)+\sum_{w:\operatorname{par}(w)=v,\,w\neq o}H(w),\qquad v\neq o,
\end{equation}
and the route count equals the sum of $H$ over root children.
The coordinate $H(o)$ is ignored; it may be set to zero for notation.
There are no travel charges in either construction below.

\section{Issue \#23: Strict-Hub Feasibility and Three-Partition}

\subsection{Source instance and target construction}

Let $m\in\N$, let $B$ be a positive integer, and let the item set be
$I=\{1,\ldots,3m\}$, equivalently $\operatorname{Fin}(3m)$ with shifted labels.
For each $i\in I$, fix an integer $a_i$ satisfying
\begin{equation}\label{eq:3p-assumptions}
 \frac{B}{4}<a_i<\frac{B}{2},\qquad
 \sum_{i\in I}a_i=mB.
\end{equation}
In particular every $a_i$ is positive. The inequalities are interpreted in $\Q$; the source weights and $B$ are integers.

\begin{definition}[Three-partition]
A three-partition is an indexed family $(P_r)_{r=1}^m$ of pairwise disjoint subsets of $I$ whose union is $I$, with
\[
 |P_r|=3,\qquad \sum_{i\in P_r}a_i=B\quad\text{for every }r.
\]
For $m=0$, the item set and family are empty, and the empty family is the unique partition in this convention.
\end{definition}
The usual nonempty problem is the restriction $m\geq1$; explicitly including $m=0$ lets the semantic statement match the existing hub theorem's zero-family boundary.

Construct a tree with distinct vertices $o$ (depot), $h$ (non-depot hub), and a leaf $\ell_i$ for every $i\in I$.
Its only edges are $\{o,h\}$ and $\{h,\ell_i\}$.
Fix
\begin{equation}\label{eq:hub-input}
 H(h)=m,\quad H(\ell_i)=1,\qquad
 s(o)=s(h)=0,\quad s(\ell_i)=a_i,\qquad T=B,
\end{equation}
and prescribe $x=E_R(H)$ and exactly $m$ routes.
Thus $x(o,h)=x(h,o)=m$, each spoke is used once in each direction, and all other flow entries vanish.
By \eqref{eq:demand}, the demands are
\begin{equation}\label{eq:hub-demands}
 d(h)=m+3m=4m,\qquad d(\ell_i)=1.
\end{equation}
The hub demand is essential: every satellite visit is followed by another hub visit, in addition to the first hub visit on each route.
The depot demand is immaterial; take $d(o)=0$.

\subsection{Exact interface to the existing hub-packing theorem}

The module \src{HubSpokeSupportExact}{\code{HubSpokeSupportExact}} uses the vertex type
\code{Node I = Option (Option I)}, with depot \code{none}, hub \code{some none},
and satellite $i$ equal to \code{some (some i)}.
Its \code{canonicalFlow m d} and \code{required m d} specialize to
\eqref{eq:hub-input}--\eqref{eq:hub-demands} when the satellite-demand function is identically $1$.
No separate root-tree adapter is needed to apply the existing hub theorem.

The reused result
\src{HubSpokePacking}{\code{HubSpokePacking.exact_iff_matrix}}
specializes exactly to
\begin{equation}\label{eq:hub-packing}
 \EF(o,x,d,m,s,B)\quad\Longleftrightarrow\quad
 \exists n\in\N^{m\times I},\quad
 \begin{cases}
  \sum_{r=1}^m n_{ri}=1,& i\in I,\\
  \sum_{i\in I}a_i n_{ri}\leq B,&r=1,\ldots,m.
 \end{cases}
\end{equation}
The general theorem permits signed service fees and budgets, but this specialization has positive item fees and a positive budget.
It also permits a row with zero satellite visits: that row is a legitimate hub-only route $[o,h,o]$.
Such rows must not be ruled out by definition; the saturation argument below rules them out for our nonempty source instances.

\begin{lemma}[Unit columns determine a partition]\label{lem:partition}
Every natural matrix $n$ in \eqref{eq:hub-packing} determines a unique assignment of each item to one route: each column has exactly one entry equal to $1$ and all other entries equal to $0$.
\end{lemma}
\begin{proof}
Every entry is at most the column sum $1$. Because the sum is $1$, at least one entry is positive, hence equal to $1$; two positive entries would make the sum at least $2$.
Set $P_r=\{i:n_{ri}=1\}$. These sets are disjoint and cover $I$.
When $I$ is empty, all these assertions about items are vacuous.
\end{proof}

\begin{lemma}[Every route saturates its budget]\label{lem:saturation}
Under \eqref{eq:3p-assumptions}, every matrix in \eqref{eq:hub-packing} satisfies
\[
 \sum_i a_i n_{ri}=B\quad\text{for every }r.
\]
\end{lemma}
\begin{proof}
Write $L_r=\sum_i a_i n_{ri}$. Each slack $B-L_r$ is nonnegative, while finite sum interchange and the unit column sums give
\[
 \sum_r L_r=\sum_i a_i\sum_r n_{ri}=\sum_i a_i=mB.
\]
Consequently $\sum_r(B-L_r)=0$, so every nonnegative slack is zero.
For $m=0$ there is no row and the conclusion is vacuous.
For $m>0$, $B>0$ then excludes every empty row.
\end{proof}

\begin{lemma}[Saturated groups contain exactly three items]\label{lem:three}
If $P\subseteq I$ and $\sum_{i\in P}a_i=B$, then $|P|=3$.
\end{lemma}
\begin{proof}
Since $B>0$, $P$ is nonempty. If $|P|\leq2$, strict upper bounds give
\[
 B=\sum_{i\in P}a_i<|P|\frac B2\leq B,
\]
a contradiction. If $|P|\geq4$, strict lower bounds give
\[
 B=\sum_{i\in P}a_i>|P|\frac B4\geq B,
\]
also a contradiction. Thus the natural number $|P|$ is $3$.
\end{proof}

\begin{theorem}[Semantic strict-hub correspondence; paper proof]\label{thm:hub}
For the construction \eqref{eq:hub-input}--\eqref{eq:hub-demands} and assumptions \eqref{eq:3p-assumptions},
\[
 \boxed{\begin{gathered}
   \EF(o,E_R(H),d,m,s,B)\\
   \Longleftrightarrow\\
   (a_i)_{i\in I}\text{ admits a three-partition}.
 \end{gathered}}
\]
\end{theorem}
\begin{proof}
For the forward direction, \eqref{eq:hub-packing} supplies a natural matrix.
Lemma \ref{lem:partition} supplies its partition $P_r$.
Lemma \ref{lem:saturation} gives each group sum $B$, and Lemma \ref{lem:three} gives each group cardinality $3$.
No arbitrary route ordering is assumed: satellite visit counts are extracted from the exact route family by the existing theorem.

For the reverse direction, given a three-partition, define $n_{ri}=1$ if $i\in P_r$ and $n_{ri}=0$ otherwise.
Disjoint coverage gives unit column sums and every row has cost exactly $B$, so \eqref{eq:hub-packing} yields exact feasibility.
Equivalently, order each triple as $(i,j,k)$ and use the full route
\[
 [o,h,\ell_i,h,\ell_j,h,\ell_k,h,o].
\]
Its four hub visits explain $d(h)=4m$; it uses the depot-hub edge once in each direction and each assigned spoke once in each direction.
It costs $a_i+a_j+a_k=B$ and has no intermediate depot visit.
The family covers every leaf exactly once, so all prescribed edge multiplicities and demands are exact.
For $m=0$, both constructions give the empty family, zero flow, and zero non-depot demands; no hub-only route is inserted.
\end{proof}

\begin{remark}[Support and exact multiplicities]
The additional existing theorem
\src{HubSpokeSupportExact}{\code{strict_hub_support_iff_exact}}
equates support-feasible and exact-feasible families for this canonical hub instance with the same demands and route count.
The statement above is deliberately made for exact flow, so it does not depend on silently replacing an arbitrary support condition by multiplicity equality.
\end{remark}

\subsection{What the correspondence establishes and what remains}

Theorem \ref{thm:hub} isolates route allocation on a simple tree after total flow is fixed.
It is a semantic equivalence for the explicitly defined inputs, reusing an existing Lean hub-packing equivalence but supplying new steps only as paper proofs.
It does not, by itself, constitute a formal polynomial-time reduction, strong NP-hardness theorem, or NP-membership proof.
It also supplies no optimality claim for a particular search upper bound.

A subsequent formalization can separate four tasks:
\begin{enumerate}
 \item encode the finite item set and positive integer hypotheses, then formalize the three elementary lemmas and the composition with \code{HubSpokePacking.exact_iff_matrix};
 \item specify the binary input/output representation, including tree, demands, flow or implicit multiplicities, weights, and fleet size;
 \item prove the size and construction-time bounds for that representation before using a complexity reduction;
 \item treat strong-hardness and NP-membership claims with their own source theorems and certificate-size hypotheses, especially when numeric counts are binary encoded.
\end{enumerate}
Only the first task's mathematical content is provided here. No further complexity classification is claimed in this package.

\section{Issue \#19: A Fixed Exterior Distinguishes Equal Subtree Totals}

\subsection{An actual tree and two-route boundary}

Use the rooted tree with vertices $o,x,h,a,b,c_1,c_2$ and edges
\[
 \{o,x\},\quad\{x,h\},\quad\{h,a\},\quad\{h,b\},\quad
 \{x,c_1\},\quad\{x,c_2\}.
\]
The internal subtree is $S=\{h,a,b\}$, separated from the exterior by edge $\{x,h\}$.
Set
\[
 s(a)=s(b)=3,\qquad s(c_1)=s(c_2)=4,\qquad
 s(o)=s(x)=s(h)=0,\qquad m=2,\quad T=7.
\]
Each of the two labeled routes crosses from $x$ to $h$ exactly once and returns from $h$ to $x$ exactly once.
Thus each route has one matched boundary excursion; neither boundary counts nor the association of the two ends with a route are being changed between the examples.

Write partial excursions as full vertex lists including their external endpoint $x$:
\begin{align*}
 Q_{ab}&=[x,h,a,h,b,h,x],& Q_0&=[x,h,x],\\
 Q_a&=[x,h,a,h,x],& Q_b&=[x,h,b,h,x].
\end{align*}
Consider the two ordered partial families
\[
 P_{\mathrm{bad}}=(Q_{ab},Q_0),\qquad
 P_{\mathrm{good}}=(Q_a,Q_b).
\]
The excursion $Q_0$ contains a hub visit and no leaf visit; it is not an empty global route, nor a forbidden intermediate return to the global depot $o$.
These excursions are partial walks in the original tree, not \code{LegalRoute h} objects with an artificially introduced depot.

Both partial families have the same aggregate directed flow:
\begin{equation}\label{eq:inside-flow}
 F(x,h)=F(h,x)=2,\quad F(h,a)=F(a,h)=1,\quad
 F(h,b)=F(b,h)=1,
\end{equation}
and zero on every other edge. They also have equal total visits in $S$: four at $h$, one at $a$, and one at $b$.
Both have total internal cost $6$, but their cost contributions per route are
\begin{equation}\label{eq:inside-cost}
 (6,0)\quad\text{and}\quad(3,3).
\end{equation}
The boundary-crossing vector of \emph{each} route is $(1,1)$ in both partial families.

\subsection{One fixed exterior and exact full routes}

For route $i$, fix the exterior walk that starts at $o$, visits $c_i$ through $x$, then attaches the subtree excursion at $x$ and finally returns to $o$.
For an excursion list $Q$ above, the resulting full route is
\[
 \mathcal E_i(Q)=[o,x,c_i]\mathbin{+\!+}Q\mathbin{+\!+}[o],
\]
where $+\!+$ is list concatenation. The endpoint $x$ is not duplicated at either join: $Q$ begins at $x$ after $c_i$ and ends at $x$ before $o$.
The exterior service contribution is always $4$.

The two completed families are explicitly
\begin{align*}
 \rho^{\mathrm{bad}}_1&=[o,x,c_1,x,h,a,h,b,h,x,o],\\
 \rho^{\mathrm{bad}}_2&=[o,x,c_2,x,h,x,o],\\
 \rho^{\mathrm{good}}_1&=[o,x,c_1,x,h,a,h,x,o],\\
 \rho^{\mathrm{good}}_2&=[o,x,c_2,x,h,b,h,x,o].
\end{align*}
Every consecutive pair is a tree edge. Every route has a nonempty task sequence and no intermediate depot visit.
For both families the nonroot parent-edge multiplicities are
\begin{equation}\label{eq:glue-h}
 H(x)=2,\qquad H(h)=2,\qquad
 H(a)=H(b)=H(c_1)=H(c_2)=1.
\end{equation}
Their aggregate directed flows are exactly $E_R(H)$, not just equal on the separator.
Formula \eqref{eq:demand}, or direct counting in the displayed lists, gives the same demands
\[
 d(x)=6,\qquad d(h)=4,\qquad d(a)=d(b)=d(c_1)=d(c_2)=1.
\]
In particular each complete route visits $x$ three times. Omitting visits caused by returning from an excursion would give an incorrect demand vector.

\begin{theorem}[Distinguishability by a fixed exterior; paper proof]\label{thm:glue}
The two partial families have equal aggregate internal flow, equal total internal cost, equal aggregate internal visits, and the same per-route boundary crossing counts and endpoint matching.
Nevertheless $P_{\mathrm{good}}$ can be attached to the fixed exterior within budget $7$, while $P_{\mathrm{bad}}$ cannot, even if its two excursions may be permuted between the exterior routes.
\end{theorem}
\begin{proof}
The asserted equalities follow by counting the four displayed excursions, as in
\eqref{eq:inside-flow}--\eqref{eq:inside-cost}.
Attaching $P_{\mathrm{good}}$ gives costs $(4+3,4+3)=(7,7)$.
Attaching $P_{\mathrm{bad}}$ gives costs $(4+6,4+0)=(10,4)$.
Either permutation of the bad excursions still attaches the contribution $6$ to an exterior contribution $4$, producing a route of cost $10>7$.
The exact full-flow and coverage conditions hold for both families by \eqref{eq:glue-h}; only the individual budget test distinguishes them.
\end{proof}

\subsection{The interface conclusion and its quantifiers}

Let $\Sigma(P)$ store the aggregate subtree flow, total internal visits and cost, and the per-route separator crossing counts with their endpoint matching.
Let $\operatorname{Glue}(P,\mathcal E)$ mean that the \emph{given} partial family $P$ can be bijectively matched to the fixed exterior slots $\mathcal E$ so that all completed routes meet their budgets.
Theorem \ref{thm:glue} proves
\begin{gather*}
 \Sigma(P_{\mathrm{good}})=\Sigma(P_{\mathrm{bad}}),\\
 \operatorname{Glue}(P_{\mathrm{good}},\mathcal E),\qquad
 \neg\operatorname{Glue}(P_{\mathrm{bad}},\mathcal E).
\end{gather*}
Consequently no function of just this summary and the exterior can exactly check gluing of \emph{each supplied partial witness}.
Permitting route relabeling does not repair the lost per-route resource contributions in this example.

The result does not exclude an existential algorithm whose state denotes all partial witnesses realizing a summary and which retains additional feasibility information or searches those witnesses.
It also does not say that a full instance is both feasible and infeasible: the common full input is feasible via $\rho^{\mathrm{good}}$.
Nor does it prove that retaining only the multiset of per-route costs suffices for arbitrary interfaces.
In the present example that multiset separates the witnesses, but general gluing may also need endpoint pairing, repeated crossings, route membership, and multiple resource contributions.

Root-branch separation has a different premise: a legal route cannot leave one root child and enter another without visiting the depot in between.
The existing \src{TreeBudgetRootFactorization}{\code{pattern_unique_root_child}} and
\code{matrix_iff_root_local} use that root-specific property.
Here $x$ and $h$ are internal vertices, so a route may visit $c_i$ and then $S$ while accumulating both service contributions.
The root factorization theorem therefore does not provide this internal interface for free.

\subsection{A bounded next specification task}
Before implementing an internal-subtree dynamic program, define a candidate interface and its gluing relation explicitly.
For a fixed outside context, the desired sound abstraction property is
\[
 \Sigma(P)=\Sigma(P')\quad\Longrightarrow\quad
 \bigl(\operatorname{Glue}(P,\mathcal E)
       \longleftrightarrow\operatorname{Glue}(P',\mathcal E)\bigr),
\]
with the admitted relabelings and matching rules stated in advance.
The current summary fails this property by Theorem \ref{thm:glue}.
A revised proposal should identify which per-route resource and boundary data it retains, prove extraction and reconstruction, and only then justify compression.
If the state instead represents an existential set of partial solutions, its invariant must say so and its transition completeness must be proved at that level.
Signed fees, when reintroduced, allow cancellation across pieces and require particular care with pruning rules based on partial costs.

\section*{Evidence, Scope, and Source Navigation}
The three-partition correspondence and gluing separation in this note are paper proofs.
No new Lean modules for Issues \#23 or \#19 are asserted.
At the initial package publication on 6 November 2026 (\code{e576911}), no LaTeX compiler was run as part of that publication or its recorded local validation.
A separately supplied six-page PDF was subsequently reviewed on the same date. This later document review does not add Lean verification, an encoding proof, or a complexity classification.
The English \code{article} source is standalone and intended for Overleaf with pdfLaTeX or XeLaTeX. This status-only source revision has not been recompiled.
Reusable source theorem links are pinned to
\href{https://github.com/huam54925-cpu/uav-routing-certification/commit/1bc992475a9bd515d77245e2bc0471e5ecab2fd5}{\code{1bc9924}}.
The exact first-package route/row contracts are discussed in the companion
\code{TREE_BUDGET_REPRESENTATION_CONTRACT.tex}.

The research tasks remain
\href{https://github.com/huam54925-cpu/uav-routing-certification/issues/23}{Issue \#23}
and \href{https://github.com/huam54925-cpu/uav-routing-certification/issues/19}{Issue \#19}.
This note establishes neither a new algorithmic technique nor a literature novelty claim.
It separates the present semantic mathematics from future formalization, encoding, and complexity work.
\end{document}
